% logs/paper/revision_M_for_C.tex
% ---------------------------------------------------------------------------------------------------
% From reviser M (mathematical sections) to reviser C (owner of paper/sec_verification.tex).
% Corrected LaTeX for the findings of referee_math_a.md and referee_math_b.md that are located in
% sec_verification.tex (and one sentence each in sec_intro.tex / the trust paragraph).  Every block
% says exactly where it goes.  Line numbers refer to paper/sec_verification.tex as it was at the start
% of the revision (= logs/paper/paper_v1/sec_verification.tex).  All labels used below exist in the
% paper (def:U-assignment, prop:maxflow, lem:hall, lem:triple, cor:pruning, lem:game-invariance,
% def:certificate, ...) or are introduced here (lem:deleted, rem:any-image).
% Nothing in this file is \input by the paper; copy the blocks.
%
% Terminology (journal referee N1), used consistently in my files from now on:
%   "certificate"        = the triple (Tbl, Lv, R_max) only;
%   "leaf condition"     = the predicates C, C_Phi, LEAF ("pairing condition C", "surplus condition
%                          C_Phi", "leaf condition LEAF");
%   "witness"            = the object that proves a single verdict (an assignment / flow behind a true
%                          verdict of LEAF, a set S of small surplus behind a false one).
% Please use the same three words in sec_verification.tex, sec_certificate.tex, app_ip.tex, sec_intro.tex.
% ---------------------------------------------------------------------------------------------------


%%%%%%%%%%%%%%%%%%%%%%%%%%%%%%%%%%%%%%%%%%%%%%%%%%%%%%%%%%%%%%%%%%%%%%%%%%%%%%%%%%%%%%%%%%%%%%%%%%%%%%
% (C1)  referee_math_a, Finding 1 (MAJOR): the leaf test of the main verifier.
%       WHERE: new unnumbered subsubsection in Section 7.3, directly AFTER the subsubsection
%       "Accelerations" (old lines 228-233) and BEFORE "Cost".
%       The lemma was checked against verify/treecheck.cpp (comment lines 628-652, validateFlow
%       lines 782-812, deletion lines 954-987) and against Definition 4.14 / Proposition 4.15.
%%%%%%%%%%%%%%%%%%%%%%%%%%%%%%%%%%%%%%%%%%%%%%%%%%%%%%%%%%%%%%%%%%%%%%%%%%%%%%%%%%%%%%%%%%%%%%%%%%%%%%

\subsubsection*{Deleting strong lines}
The main verifier uses one further device, and this one changes the assignments on which its true
verdicts rest. Let $\theta\ge1$. As soon as a line has been found strong it is \emph{deleted}: it gives
up its two cells and is ignored in all later searches, for single lines, pairs and triples alike. To
justify this, let $D\subseteq\live{B}$ be a set of lines and $L_D=\live{B}\setminus D$. For
$T\subseteq L_D$, a \emph{$T$-assignment relative to $D$} is a map $\alpha$ from a set of cells to $L_D$
such that $c\in\alpha(c)$ and $c\notin F$ for every assigned cell $c$, every line of $L_D\setminus T$
holds exactly two cells and every line of $T$ holds at least two; its \emph{excess} is
$\sum_{l\in T}(|\alpha^{-1}(l)|-2)$. The lines of $D$ hold no cell. Put
\[
  \mu_D(T)=\min\{\sur_F(S) : T\subseteq S\subseteq L_D\}.
\]
For $D=\emptyset$ these are the notions of \cref{def:U-assignment} and the number $\mu(T)$.

\begin{lemma}[Deleted lines]\label{lem:deleted}
Let $\theta\ge 1$, let $D\subseteq\live{B}$ be a set of lines with $\mu(\{d\})\ge\theta$ for every
$d\in D$, and let $T\subseteq L_D$.
\begin{enumerate}[label=\textup{(\alph*)}]
\item If there is a $T$-assignment relative to $D$ of excess at least $\theta$, then $\mu(T)\ge\theta$.
\item $\min(\theta,\mu(T))=\min(\theta,\mu_D(T))$. If $\pi$ is a pairing for $(B,F)$, then $\mu_D(T)$ is
      the number of augmentations along augmenting paths from lines of $T$ that are performed when one
      starts from the restriction of $\pi$ to the cells held by the lines of $L_D$, uses lines of $L_D$
      only, and continues until there is no such path.
\end{enumerate}
\end{lemma}

\begin{proof}
(a) Let $T\subseteq S\subseteq\live{B}$; we have to show $\sur_F(S)\ge\theta$. If $S$ contains a line
$d\in D$, then $\sur_F(S)\ge\mu(\{d\})\ge\theta$ by the definition of $\mu$. Otherwise $S\subseteq L_D$.
The cells held by the lines of $S$ are distinct cells of $\N(S)\setminus F$, and their number is at least
$2|S|+\theta$, because $T\subseteq S$, the lines of $S\setminus T$ hold two cells each and the lines of
$T$ hold at least $2|T|+\theta$ cells together. Hence $|\N(S)\setminus F|\ge 2|S|+\theta$.

(b) By definition, $\mu(T)$ is the smaller of $\mu_D(T)$ and of the minimum of $\sur_F(S)$ over the sets
$S\supseteq T$ of live lines that contain a line of $D$; the latter is at least $\theta$, as in (a). This
gives the first assertion. \Cref{def:U-assignment}, \cref{prop:maxflow} and its proof use no property of
the set $\live{B}$ except that $\pi$ assigns two cells of $l\setminus F$ to each of its lines $l$, all
these cells being distinct; they apply verbatim to the set $L_D$ of lines and to the restriction of
$\pi$. Part (d) of the proposition, read for $L_D$, is the second assertion.
\end{proof}

By induction on the order of deletion, every deleted line is indeed strong: when a line $d$ is tested,
the lines deleted before it are strong by the induction hypothesis, and if the test finds a
$\{d\}$-assignment of excess $\theta$ relative to them, then $\mu(\{d\})\ge\theta$ by part (a). Hence
\cref{lem:deleted} applies to every set $D$ of deleted lines that occurs, and single lines, pairs and
triples may be tested in the network without the deleted lines: by (a) a true verdict of such a test is
correct, and by (b) so is a false one. For a verifier that deletes strong lines the properties to be
validated behind a true verdict for a set $T$ are therefore: no cell is assigned twice; every assigned
cell lies on its line and outside $F$; the deleted lines hold no cell; every other live line outside $T$
holds exactly two cells; and the lines of $T$ hold $2|T|+\theta$ cells together. For $\theta=1$ a line
$l$ is strong if and only if there is an augmenting path from $l$ with respect to the pairing
(\cref{prop:maxflow}\,(d)); the main verifier finds all these lines by a single backward search from the
unassigned cells and, with the option \fn{--paranoid} that was used in the certifying run, constructs
and validates the assignment for each of them nevertheless.


%%%%%%%%%%%%%%%%%%%%%%%%%%%%%%%%%%%%%%%%%%%%%%%%%%%%%%%%%%%%%%%%%%%%%%%%%%%%%%%%%%%%%%%%%%%%%%%%%%%%%%
% (C2)  referee_math_a, Finding 1 (i) and symptom (b): Remark 7.3 (rem:true-cert), old lines 213-226.
%       REPLACE the whole remark by the following.  Changes: the pair rule is covered; the main
%       verifier is said to validate the modified properties of (C1); title uses "witness" (N1).
%       (Keep the label rem:true-cert, it is cited from Table 7.)
%%%%%%%%%%%%%%%%%%%%%%%%%%%%%%%%%%%%%%%%%%%%%%%%%%%%%%%%%%%%%%%%%%%%%%%%%%%%%%%%%%%%%%%%%%%%%%%%%%%%%%

\begin{remark}[True verdicts carry explicit witnesses]\label{rem:true-cert}
For the soundness of a verifier only the implication ``the algorithm returns true $\Rightarrow\LEAF$''
matters, and this implication does not need the max-flow min-cut theorem. Suppose
$\textsc{Gain}(\mathit{own},U,\theta)=\theta$, and let $\mathit{own}'$ be the final assignment: no cell is
assigned twice, every live line outside $U$ holds exactly two cells, and the lines of $U$ hold $2|U|+\theta$
cells together. For every set $S\supseteq U$ of live lines, the lines of $S$ then hold $2|S|+\theta$
distinct cells of $\N(S)\setminus F$, so $\sur_F(S)\ge\theta$ (\cref{prop:maxflow}\,(a)). Every $F$-connected
set $S$ with $|S|\ge3$ contains an $F$-connected triple $T$ (\cref{lem:triple}), and $T$ either has
$\textsc{Gain}=\theta$ itself, or contains a strong line $l$, or, if the second pruning rule below is
used, contains a pair $\{a,b\}$ with $\textsc{Gain}=\theta$; in the last two cases the same counting
applies with $U=\{l\}$ or $U=\{a,b\}$. An implementation can therefore validate each true verdict
literally, by checking the final assignments against these three properties. The main verifier does
so, in the modified form described under ``Deleting strong lines'' below (\cref{lem:deleted}). False
verdicts affect only completeness: a wrong ``false'' can make a verifier reject a valid certificate,
never accept an invalid one.
\end{remark}


%%%%%%%%%%%%%%%%%%%%%%%%%%%%%%%%%%%%%%%%%%%%%%%%%%%%%%%%%%%%%%%%%%%%%%%%%%%%%%%%%%%%%%%%%%%%%%%%%%%%%%
% (C3)  referee_math_a, Finding 1 (ii): Table 7 (tab:verifiers), first row, middle column,
%       old lines 259-261.  REPLACE
%          "\cref{alg:leaf} with both pruning rules, every flow behind a true verdict validated
%           literally (\cref{rem:true-cert})"
%       BY
%%%%%%%%%%%%%%%%%%%%%%%%%%%%%%%%%%%%%%%%%%%%%%%%%%%%%%%%%%%%%%%%%%%%%%%%%%%%%%%%%%%%%%%%%%%%%%%%%%%%%%

%   \cref{alg:leaf} with both pruning rules and deletion of strong lines (\cref{lem:deleted}); for
%   $\theta=1$ the strong lines are found by one backward search; every assignment behind a true
%   verdict validated literally (\cref{rem:true-cert})

% In "Accelerations" (old lines 232-233) the sentence "None of this is needed for correctness, and
% independent implementations differ in what they use" stays true for the three devices named there;
% add after it:  "The main verifier uses a fourth device, which needs an argument of its own."


%%%%%%%%%%%%%%%%%%%%%%%%%%%%%%%%%%%%%%%%%%%%%%%%%%%%%%%%%%%%%%%%%%%%%%%%%%%%%%%%%%%%%%%%%%%%%%%%%%%%%%
% (C4)  referee_math_a, Finding 1 (iii) and symptom (a): "Cost", old lines 239-241.
%       The number 869,185 is NOT the number of single-line gains of Algorithm 2 (that would be about
%       four million: one per live line of each of the 41,882 leaves with theta >= 1).  According to the
%       source (treecheck.cpp, counter nSingle, lines 954-975) it counts only the single-line flows for
%       theta >= 2, from lines that have at least theta+2 cells outside F and from which an augmenting
%       path exists at all.  Please re-read the three counters from results/final/treecheck.json
%       (leaf_triple_flows = 7391 is there; take the names of the other two from the file) and REPLACE
%          "over all $106{,}723$ leaves the main verifier computes $869{,}185$ single-line flows,
%           $77{,}268$ flows for pairs and only $7{,}391$ flows for triples."
%       BY (numbers to be confirmed from the json by reviser C):
%%%%%%%%%%%%%%%%%%%%%%%%%%%%%%%%%%%%%%%%%%%%%%%%%%%%%%%%%%%%%%%%%%%%%%%%%%%%%%%%%%%%%%%%%%%%%%%%%%%%%%

%   over all $106{,}723$ leaves the main verifier, which deletes strong lines, computes $869{,}185$
%   single-line flows (these are counted for $\theta\ge2$ only, and only for lines that have at least
%   $\theta+2$ cells outside $F$ and from which an augmenting path exists at all; for $\theta=1$ the
%   strong lines are found by the backward search), $77{,}268$ flows for pairs and only $7{,}391$ flows
%   for triples.


%%%%%%%%%%%%%%%%%%%%%%%%%%%%%%%%%%%%%%%%%%%%%%%%%%%%%%%%%%%%%%%%%%%%%%%%%%%%%%%%%%%%%%%%%%%%%%%%%%%%%%
% (C5)  referee_math_a, Finding 2: proof of Proposition 7.2 (prop:leaf-correct), old lines 200-203.
%       REPLACE the two sentences
%          "The loop of lines \ref{ln:pair1}--\ref{ln:pair2} is the augmenting-path algorithm for the
%           network of \cref{lem:hall}, in which every live line has capacity two; it ends with a
%           pairing if one exists and fails otherwise, which decides (L2)."
%       BY the following (the search is restricted to paths that start at the one line being
%       processed, so "fails otherwise" needs its own argument; this one avoids max-flow min-cut):
%%%%%%%%%%%%%%%%%%%%%%%%%%%%%%%%%%%%%%%%%%%%%%%%%%%%%%%%%%%%%%%%%%%%%%%%%%%%%%%%%%%%%%%%%%%%%%%%%%%%%%

The loop of lines \ref{ln:pair1}--\ref{ln:pair2} decides (L2). Throughout the loop every cell is held
by at most one line, a line holds only cells of its own outside $F$, and every line holds at most two
cells: an augmentation from $l$ gives one more cell to $l$, which held at most one, and leaves the
number of cells of every other line unchanged. If the loop ends, every live line holds exactly two
cells, and $\mathit{own}$ is a pairing for $(B,F)$. If \textsc{Augment}$(\mathit{own},\{l\})$ fails in
line~\ref{ln:pair2}, let $R$ be the set of lines that are reachable from $l$ by alternating paths
(\cref{def:U-assignment}; the definition of alternating paths and fact~(iii) after it do not require
that the lines hold two cells). Every cell of $\N(R)\setminus F$ is held by a line of $R$: a cell
$c\in l'\setminus F$ with $l'\in R$ that is not assigned would complete an augmenting path from $l$,
and a cell that is held by a line $l''\ne l'$ makes $l''$ reachable. At this moment every line of $R$
holds at most two cells and $l\in R$ holds at most one; hence $|\N(R)\setminus F|\le 2|R|-1$, and there
is no pairing for $(B,F)$ by \cref{lem:hall}\,(b).


%%%%%%%%%%%%%%%%%%%%%%%%%%%%%%%%%%%%%%%%%%%%%%%%%%%%%%%%%%%%%%%%%%%%%%%%%%%%%%%%%%%%%%%%%%%%%%%%%%%%%%
% (C6)  referee_math_b, Finding 3: Section 7.1, old lines 40-42, the unproved side remark
%          "For soundness it is not even necessary that the representative stored for a state is the
%           \emph{minimal} image; it is enough that it is the image of the state under one element of
%           $\G$. Minimality matters only for finding the entries of the two files."
%       Proposition 7.1 does not cover it (its Seen = V consists of canonical forms).
%       REPLACE these sentences in Section 7.1 by the pointer
%          "A verifier that stores an arbitrary image of a state instead of the minimal one is
%           still sound, although \cref{prop:verify-correct} does not apply to it literally; see
%           \cref{rem:any-image}."
%       and INSERT the following remark directly AFTER the proof of Proposition 7.1
%       (prop:verify-correct, after old line 137).
%%%%%%%%%%%%%%%%%%%%%%%%%%%%%%%%%%%%%%%%%%%%%%%%%%%%%%%%%%%%%%%%%%%%%%%%%%%%%%%%%%%%%%%%%%%%%%%%%%%%%%

\begin{remark}[Representatives that are not minimal]\label{rem:any-image}
Suppose that in \cref{alg:verify} the function $\canon$ is replaced by any function $r$ that maps every
state to the image of this state under some element of $\G$, which may depend on the state, and that
the algorithm returns \accept. Then the root is still \WON. Indeed, every state $s$ that was processed
is \WON, by induction on the rank of $s$ (\cref{lem:rank}; a state and its images under $\G$ have the
same rank, since the rank depends only on $t$, $|F|$ and the kind of the state). If $s$ is a P-state
in $\Lv$, then $\LEAF$ was required. If $s$ is a P-state outside $\Lv$, then $t<\Rmax$ and the
falsity of $\LOST$ were required, and for every child $s'$ of $s$ of type (a), (c) or (d) that is not
\VOID{} the state $r(s')$ was processed and has smaller rank than $s$; so $r(s')$ is \WON, and so is
$s'$ by \cref{lem:game-invariance}\,(c). If $s$ is a Q-state, the cell $m=\Tbl(s)\notin B\cup F$ that
was required is a table move by the same argument, applied to the child (i) and to the child (ii) if
it exists and is not \VOID. Finally the root is processed, because it is fixed by every element of
$\G$. In this situation the triple $(\Tbl,\Lv,\Rmax)$ need not be a valid certificate in the sense of
\cref{def:certificate}, whose reachable set consists of canonical forms; minimality of the stored
representative matters for finding the entries of the two files, which are listed by canonical form.
\end{remark}

% NOTE for (C6): please check the label of the rank lemma in sec_game.tex (Lemma 5.4); if it is not
% lem:rank, use the right one.  [M will confirm the label in logs/paper/response_M.md.]


%%%%%%%%%%%%%%%%%%%%%%%%%%%%%%%%%%%%%%%%%%%%%%%%%%%%%%%%%%%%%%%%%%%%%%%%%%%%%%%%%%%%%%%%%%%%%%%%%%%%%%
% (C7)  referee_math_b, Finding 5, first bullet: Algorithm 1 (old lines 69-73) and the comment
%       "The children" (old lines 100-103).  In Definition 5.2, Lemma 5.12, Remarks 5.13, 5.18,
%       Proposition C.3 and Example C.5 the letters a', b' are LINES; in Algorithm 1 they are cells.
%       Rename the two cells to p, q.  REPLACE old lines 70-71 by
%%%%%%%%%%%%%%%%%%%%%%%%%%%%%%%%%%%%%%%%%%%%%%%%%%%%%%%%%%%%%%%%%%%%%%%%%%%%%%%%%%%%%%%%%%%%%%%%%%%%%%

%      \ForEach{$p\in a\setminus(F\cup\{c\})$ and $q\in b\setminus(F\cup\{c\})$}{
%        $Q_a\gets a\setminus\{c,p\}$;\quad $Q_b\gets b\setminus\{c,q\}$\;

%       and in the comment "The children" REPLACE
%          "$Q_a$ is determined by the one cell $a'$ that is left out, and the condition
%           $F\cap a\subseteq Q_a$ says $a'\notin F$."
%       BY
%          "$Q_a$ is determined by the one cell $p$ that is left out, and the condition
%           $F\cap a\subseteq Q_a$ says $p\notin F$; likewise for $Q_b$ and $q$."
%       (In my files the letters a', b' now denote lines only, and the board is never written next to
%       the sets Q_a, Q_b or the Q-states; see response_M.md.)


%%%%%%%%%%%%%%%%%%%%%%%%%%%%%%%%%%%%%%%%%%%%%%%%%%%%%%%%%%%%%%%%%%%%%%%%%%%%%%%%%%%%%%%%%%%%%%%%%%%%%%
% (C8)  referee_math_a, Finding 4: Section 7.6 "The integer program", old lines 400-401.
%       "weaker" is the opposite of the terminology of Sections 4 and 9 (LEAF implies C, so LEAF is
%       the STRONGER condition), and "when F \neq \emptyset" wrongly suggests equivalence for F empty
%       (true only for k <= 10, Proposition 4.20 / Example 4.11).  REPLACE
%          "Among the $656$ instances that hold there are $160$ for which $\LEAF$ is false: $\LEAF$ is
%           strictly weaker than $\CC$ when $F\neq\emptyset$."
%       BY
%%%%%%%%%%%%%%%%%%%%%%%%%%%%%%%%%%%%%%%%%%%%%%%%%%%%%%%%%%%%%%%%%%%%%%%%%%%%%%%%%%%%%%%%%%%%%%%%%%%%%%

Among the $656$ instances that hold there are $160$ for which $\LEAF$ is false: $\LEAF$ is a strictly
stronger condition than $\CC$ (it fails on instances on which $\CC$ holds), not only in the situation
of \cref{ex:C-not-leaf}.


%%%%%%%%%%%%%%%%%%%%%%%%%%%%%%%%%%%%%%%%%%%%%%%%%%%%%%%%%%%%%%%%%%%%%%%%%%%%%%%%%%%%%%%%%%%%%%%%%%%%%%
% (C9)  referee_math_a, Finding 1, "trust" sentences: Section 7.7 (old lines 459-463) and
%       sec_intro.tex (old lines 189-193, 198-199).  "at least one implementation of these two
%       algorithms" is inaccurate for the main verifier, whose leaf test is Algorithm 2 with the
%       deletion of strong lines.  Suggested wording for the second and third item of Section 7.7:
%%%%%%%%%%%%%%%%%%%%%%%%%%%%%%%%%%%%%%%%%%%%%%%%%%%%%%%%%%%%%%%%%%%%%%%%%%%%%%%%%%%%%%%%%%%%%%%%%%%%%%

%  \item that the definition of validity used by the verifiers (\cref{ssec:what}) is the one of
%    \cref{def:certificate}, and that \cref{alg:verify,alg:leaf} are correct
%    (\cref{prop:verify-correct,prop:leaf-correct}), as is the variant of \cref{alg:leaf} with deleted
%    lines that the main verifier uses (\cref{lem:deleted});
%  \item at least one implementation of these algorithms (\fn{verify/spotcheck.py} implements
%    \cref{alg:leaf} as printed, \fn{verify/treecheck.cpp} the variant with deleted lines), together
%    with ...
%
%  and in sec_intro.tex, after "... and proves them correct
%  (\cref{prop:verify-correct,prop:leaf-correct})" add:
%  "; the main verifier uses a variant of the second algorithm, which is justified in
%  \cref{lem:deleted}".


%%%%%%%%%%%%%%%%%%%%%%%%%%%%%%%%%%%%%%%%%%%%%%%%%%%%%%%%%%%%%%%%%%%%%%%%%%%%%%%%%%%%%%%%%%%%%%%%%%%%%%
% (C10) Labels and wording in MY files that sec_verification.tex / sec_certificate.tex / app_*.tex
%       cite.  No label of Sections 2-5, Appendix C or Appendix D is renamed or removed by reviser M.
%       New labels will be listed at the end of logs/paper/response_M.md (section "For reviser C").
%%%%%%%%%%%%%%%%%%%%%%%%%%%%%%%%%%%%%%%%%%%%%%%%%%%%%%%%%%%%%%%%%%%%%%%%%%%%%%%%%%%%%%%%%%%%%%%%%%%%%%


%%%%%%%%%%%%%%%%%%%%%%%%%%%%%%%%%%%%%%%%%%%%%%%%%%%%%%%%%%%%%%%%%%%%%%%%%%%%%%%%%%%%%%%%%%%%%%%%%%%%%%
% ADDED 17:40 UTC (second batch; everything above is unchanged)
%%%%%%%%%%%%%%%%%%%%%%%%%%%%%%%%%%%%%%%%%%%%%%%%%%%%%%%%%%%%%%%%%%%%%%%%%%%%%%%%%%%%%%%%%%%%%%%%%%%%%%

%%%%%%%%%%%%%%%%%%%%%%%%%%%%%%%%%%%%%%%%%%%%%%%%%%%%%%%%%%%%%%%%%%%%%%%%%%%%%%%%%%%%%%%%%%%%%%%%%%%%%%
% (C11) referee_math_a, Finding 5, second location: sec_discussion.tex, paragraph "LEAF versus C"
%       (old lines 261-263; this paragraph is yours, I only edit Section 9.1 and Proposition 9.5).
%       For theta = 0 the two conditions agree only if F is B-admissible (otherwise C holds vacuously
%       and LEAF fails by (L1)).  REPLACE
%          "The two conditions agree when all stones are revealed and when none is revealed and
%           $t\le 10$"
%       BY
%          "The two conditions agree, for admissible $F$, when all stones are revealed, and when
%           none is revealed and $t\le 10$"
%%%%%%%%%%%%%%%%%%%%%%%%%%%%%%%%%%%%%%%%%%%%%%%%%%%%%%%%%%%%%%%%%%%%%%%%%%%%%%%%%%%%%%%%%%%%%%%%%%%%%%

%%%%%%%%%%%%%%%%%%%%%%%%%%%%%%%%%%%%%%%%%%%%%%%%%%%%%%%%%%%%%%%%%%%%%%%%%%%%%%%%%%%%%%%%%%%%%%%%%%%%%%
% (C12) Journal referee T2 / E6: Table 2 (tab:notation, in sec_board.tex) is now extended by the
%       notation of Section 5 and cited from Sections 2 and 5.  Please cite it once in Section 1.8
%       (sec_intro.tex, "Organization"):  "The notation is collected in \cref{tab:notation}."
%%%%%%%%%%%%%%%%%%%%%%%%%%%%%%%%%%%%%%%%%%%%%%%%%%%%%%%%%%%%%%%%%%%%%%%%%%%%%%%%%%%%%%%%%%%%%%%%%%%%%%

%%%%%%%%%%%%%%%%%%%%%%%%%%%%%%%%%%%%%%%%%%%%%%%%%%%%%%%%%%%%%%%%%%%%%%%%%%%%%%%%%%%%%%%%%%%%%%%%%%%%%%
% (C13) Replication referee, finding 7 (counting conventions): Section 5.2 used to say "Different
%       data may define the same child".  For a parent that is NOT void this is false, and Section 5.2
%       now proves it:  \cref{lem:children-distinct}  ("Children are counted once": for a P-state that
%       is not VOID, different data (l,e) resp. ({a,b},Q_a,Q_b) define different children, and a child
%       (c) is never a child (d); the two children of a Q-state are of different kinds).
%       So the two wordings of the caption of Table 3 ("counted every time it is generated") and of
%       Appendix A.5 ("for every pair of an expanded state and one of its non-void children (the
%       children of one state are pairwise distinct)") describe the same number.  Please use ONE
%       wording in both places and cite the lemma, e.g.
%          "A child is counted once for every expanded state of which it is a child; by
%           \cref{lem:children-distinct} this is the number of data $(l,e)$, respectively
%           $(\{a,b\},Q_a,Q_b)$, that are generated in lines \ref{ln:childa}--\ref{ln:childd} of
%           \cref{alg:verify}.  Children of different states that have the same canonical form are
%           counted separately."
%       Related (replication finding 6): Algorithm 1, line 25 (Push) drops any VOID child silently,
%       although a child (a) or (i) of a state that is not VOID is never VOID (\cref{lem:adm-closed});
%       this is stated in Section 5.2 after \cref{def:states}.
%%%%%%%%%%%%%%%%%%%%%%%%%%%%%%%%%%%%%%%%%%%%%%%%%%%%%%%%%%%%%%%%%%%%%%%%%%%%%%%%%%%%%%%%%%%%%%%%%%%%%%

%%%%%%%%%%%%%%%%%%%%%%%%%%%%%%%%%%%%%%%%%%%%%%%%%%%%%%%%%%%%%%%%%%%%%%%%%%%%%%%%%%%%%%%%%%%%%%%%%%%%%%
% (C14) Journal referee N4 (the letter Q).  The symbol of the BOARD is changed from Q to \Omega by
%       one line appended to macros.tex  (\renewcommand{\Q}{\Omega}); the 7x7x7 board \Qseven becomes
%       \Omega_7 automatically.  Nothing to do in your files as long as the board is written \Q
%       (it is, everywhere).  Please never write a literal $Q$ for the board.  The letter Q then
%       means: Q-states \Qst(B,F,t), and the stone sets Q_a, Q_b of the danger children (kept,
%       because they occur in Algorithm 1 and in Appendix A; in my files they are always called
%       "the revealed stones Q_a on a").  The letters a', b' are lines only (C7).
%
% (C15) Journal referee N1 in the remaining places of YOUR files that speak about Sections 3-5:
%       "A pairing is a certificate that is easy to check" (Section 3.4) is now "a witness";
%       Section 4 is now called "Leaf conditions"; Definition 4.4 "The pairing condition C",
%       Definition 4.7 "The surplus condition C_Phi and the leaf condition LEAF".
%       \cref{...} to these labels is unaffected (labels unchanged: sec:leaf, def:C, def:leaf, ...).
%       Wording to use elsewhere:  "the pairing condition $\CC(B,F,k)$", "the surplus condition
%       $\CPhi$", "the leaf condition $\LEAF$" (not: leaf certificate, leaf predicate, pairing
%       certificate).  "Predicate" is still used for $\LOST$.
%
% (C16) Journal referee P9 (Zermelo): in sec_game.tex and in Section 9.1 the wording is now
%       "by the determinacy of finite games with perfect information, usually attributed to
%       Zermelo~\cite{Zermelo1913} and proved by backward induction on the game tree".
%       If sec_intro.tex (Section 1.6) mentions Zermelo, please use the same wording.
%%%%%%%%%%%%%%%%%%%%%%%%%%%%%%%%%%%%%%%%%%%%%%%%%%%%%%%%%%%%%%%%%%%%%%%%%%%%%%%%%%%%%%%%%%%%%%%%%%%%%%


%%%%%%%%%%%%%%%%%%%%%%%%%%%%%%%%%%%%%%%%%%%%%%%%%%%%%%%%%%%%%%%%%%%%%%%%%%%%%%%%%%%%%%%%%%%%%%%%%%%%%%
% ADDED at the end of reviser M's phase (third batch; information only, nothing above is changed).
% Thank you for integrating (C1)-(C7), (C12), (C13); I have read the result in sec_verification.tex
% (lem:deleted, rem:any-image, "The leaf test of the main verifier") and have no objection.
%%%%%%%%%%%%%%%%%%%%%%%%%%%%%%%%%%%%%%%%%%%%%%%%%%%%%%%%%%%%%%%%%%%%%%%%%%%%%%%%%%%%%%%%%%%%%%%%%%%%%%
%
% (C17) New labels in my files (all labels that existed before are unchanged):
%         lem:children-distinct   sec_game.tex, Section 5.2  "Different data define different children"
%         ex:first-rounds         sec_game.tex, Section 5.2  worked example (root, child (a), an occupied
%                                 cell, a threat child, a danger child); it cites \cref{tab:opening},
%                                 \cref{ssec:walk}, \cref{ex:T8}, \cref{fact:H} -- please keep these labels
%         fig:patterns            sec_game.tex, Section 5.3  figure of the patterns E2 and E3
%       The numbering inside Section 5 has shifted (old Definition 5.2 = new 5.4, old Lemma 5.7 = new
%       5.9, old Theorem 5.14 = new 5.16, ...).  Everything in the paper uses \cref; if README.md or
%       another archive document quotes statement numbers of Section 5, they have to be updated after
%       the final build.
%
% (C18) "strong": the definition (a line l with mu({l}) >= theta, a pair with mu({a,b}) >= theta) is now
%       also printed in Section 4.3 directly after \cref{cor:pruning}, so Section 7.3 may refer to it.
%
% (C19) Proposition 4.22 (prop:busy10) is now stated for |B| <= 17 (referee A, finding 7; the label is
%       unchanged).  Section 6.4 quotes it for sets of ten busy cells, which remains correct.
%
% (C20) Files that are cited in MY sections and therefore have to be in the archive (D4); all exist in
%       the project tree:
%         logs/canon/autgroup.cpp, logs/canon/autgroup.log            (Computed Fact 2.9)
%         logs/minstones/run.log, logs/core/test_core.log             (Computed Facts 3.10, 4.23)
%         logs/paper/secA_check1.py, logs/paper/secA_checks.log       (Computed Facts 3.10, 4.23)
%         results/regression/ip_regression.json                       (Computed Fact 4.23)
%         logs/proof/fuzz_cprime.py                                   (end of Appendix C.1)
%         logs/proof/e23min.cpp, logs/proof/e23min.log                (Computed Fact C.4)
%         scripts/check_7cube_midplanes.py                            (Remark 9.4 (ii))
%       scripts/check_7cube_midplanes.py now READS paper/app_seven.tex (it checks the printed table of
%       Appendix D), so the archive must contain paper/app_seven.tex at that relative path; its output
%       is logs/paper/check_7cube_midplanes.out (regenerated, exit code 0).
%       My own check scripts for the revision (not cited in the paper) are in logs/paper/revM_checks/.
%
% (C21) Not done by me, for your decision:
%       * replication referee, finding 3: test vectors for LOST with a positive answer (Appendix A.6);
%         I only added the sentence on "x in F" after Definition 5.4 (old 5.2);
%       * journal referee E5/E10 and the table of Section 6 of his report (moving Section 4.3, the
%         five-step list at the end of 4.3, the last page of Section 5.10, Appendix C): these are
%         moves between files and belong to your second phase (D3); I moved nothing;
%       * sec_discussion.tex, "The shape of the certificate": "from round $8$ on he can have a danger
%         pair" still has a pronoun (your part of the file); likewise (C11).
